\subsection{A sharp energy-and-bath threshold for the contrast family}\label{sec:energythreshold}
The contrast theorem also determines exactly which energy/bath choices
exclude ordinary cooling \emph{uniformly over the entire rate family}.
Shift $E_1=0$, let $a=E_2>0$, and let
$\mu=n^{-1}\sum_iE_i$ be the arithmetic mean energy, not the bath mean.

\begin{theorem}[Uniform cooling threshold]\label{thm:energythreshold}
Fix distinct ordered energies and a positive finite bath. The following are
equivalent:
\begin{enumerate}
\item Every positive reversible generator satisfying
$k_{ik}\leq k_{jk}$ for $i<j<k$ has no ordinary asymptotic cooling effect
in weighted Euclidean distance.
\item Either $\mu\leq a$, or
\begin{equation}\label{eq:energythreshold}
\beta_b\leq\beta_*:=\frac1a\log\frac{\mu}{\mu-a}.
\end{equation}
\end{enumerate}
Under condition 2, absence is stronger: it holds at every time in weighted
Euclidean distance, for asymptotic relative entropy at any slow multiplicity,
and asymptotically in every fixed norm when the slow mode is simple. The
same universal equivalence holds for the larger family in which the
contrast matrix is Metzler. Failure of condition 2 does not by itself assert
existence under the additional original $F$ inequalities.
\end{theorem}
\begin{proof}
For cooling, each contrast amplitude $R_k=-g_k$ is strictly positive.
With $\alpha_k=Q_k/q_{k+1}=\sum_{i\leq k}e^{\beta_b(E_{k+1}-E_i)}$,
\[
R_k=\frac{e^{-\beta E_{k+1}}}
{\sqrt{\eta_k}Z}\left(\alpha_k-
\sum_{i=1}^ke^{\beta(E_{k+1}-E_i)}\right),\qquad \beta<\beta_b.
\]
The logarithmic derivative of the bracket is negative and strictly
decreasing, while $(\log Z)''=\operatorname{Var}_{\pi(\beta)}E>0$.
Therefore $(\log R_k)''<0$: every cooling contrast amplitude is strictly
log-concave. It follows that $g_k'(0)\geq0$ implies $g_k'(\beta)>0$ for
every $0<\beta<\beta_b$, including the equality case at $0$.

Direct differentiation gives
\begin{equation}\label{eq:contrastD0}
n^2\sqrt{\eta_k}\,g_k'(0)=
n\sum_{i=1}^k\left[
e^{\beta_b(E_{k+1}-E_i)}(E_{k+1}-\mu)+(\mu-E_i)\right].
\end{equation}
For $k=1$ this has the sign of
$\mu-e^{\beta_ba}(\mu-a)$. It is nonnegative precisely under condition 2.
We show that this one condition makes every other expression nonnegative.
Put $y=E_{k+1}$. If $y\geq\mu$, the sum equals
$\alpha_k(y-\mu)+k\mu-\sum_{i\leq k}E_i>0$, because the mean of a proper
ordered prefix is below the full arithmetic mean.
If $y<\mu$, necessarily $a<\mu$. For every $i\leq k$,
\[
\frac{\log(\mu-E_i)-\log(\mu-y)}{y-E_i}
=\frac1{y-E_i}\int_{E_i}^{y}\frac{dt}{\mu-t}
\geq\frac1a\int_0^a\frac{dt}{\mu-t}=\beta_*.
\]
The inequality holds because the integrand increases: for $i=1$, the
interval starts at zero and extends at least to $a$; for $i>1$ its entire
interval starts at or above $a$. Thus $\beta_b\leq\beta_*$ implies
$e^{\beta_b(y-E_i)}\leq(\mu-E_i)/(\mu-y)$, making each summand in
\eqref{eq:contrastD0} nonnegative. Consequently all $g_k$ strictly increase
throughout cooling. Their negatives are ordered componentwise for two
cooling preparations, and the positive contrast semigroup gives the same
finite-time norm ordering as in Theorem~\ref{thm:orderedk}. The spectral
projection argument gives the asymptotic conclusions.

Conversely, suppose $a<\mu$ and $\beta_b>\beta_*$. The first contrast has
$g_1'(0)<0$, so $|g_1|$ initially increases with inverse temperature and
gives finite positive cooling witnesses. Realize $h^{(1)}$ as the simple
slow mode with the established rank-one construction
\[
M=-I+q\mathbf1^\top+\epsilon
\frac{D h^{(1)}(h^{(1)})^\top}{\eta_1}
\]
for sufficiently small $0<\epsilon<1$. Positivity and the simple slow
eigenvalue follow as before. Its reduced conductances are
$k_{ij}=1+\epsilon h_i^{(1)}h_j^{(1)}/\eta_1$. For every $i<j<k$,
$h_k^{(1)}=0$, so $k_{ik}=k_{jk}=1$. Thus all $G$ inequalities hold
at equality. Its contrast matrix is diagonal, so it also belongs to the
larger family. This proves sharpness of the uniform threshold.
\end{proof}

For four equally spaced energies $(0,a,2a,3a)$, the threshold is
$\beta_ba\leq\log3$. Throughout that range, failure of all triplet tests
really does exclude ordinary cooling as well as the other effects covered
above. Outside that range, the larger $G$-only family permits cooling, but
the original $F$-constrained equal-spacing question is not settled by the
sharpness example. For three equally spaced energies, $\mu=a$, so the
$G$-only family has no cooling effect at any finite positive bath temperature.

\begin{corollary}[A spacing-independent energy-span exclusion]\label{cor:span}
For $n\geq3$, write $c=E_n-E_1$. If
\[
\beta_b c\leq \frac{n}{n-2},
\]
every positive reversible generator with a Metzler contrast matrix has no
finite-time cooling reversal in weighted Euclidean distance. The asymptotic
conclusions of Theorem~\ref{thm:energythreshold} hold as well. In particular,
these exclusions apply under failure of all original triplet predicates.
The constant $n/(n-2)$ is optimal uniformly over energy spacings in the
larger $G$-only family; optimality under the additional $F$ constraints is
not asserted.
\end{corollary}
\begin{proof}
Shift $E_1=0$, set $r=a/c$, and put $d=n-2$. The case $\mu\leq a$ is
already covered. Otherwise $0<r<d/(d+1)$ and
$\mu\leq c(d+r)/n$. The threshold decreases as $\mu$ increases, so
\[
c\beta_*\geq
\frac{\log(1+r/d)-\log(1-(d+1)r/d)}{r}
>\frac nd.
\]
The strict inequality follows by expanding the logarithms: the constant
term is $n/d$, and the coefficient of $r^{k-1}$ for every $k\geq2$ is
$[(-1)^{k+1}+(d+1)^k]/(k d^k)>0$. Both series converge on the stated
interval. This proves absence, including equality at the advertised cutoff.

For sharpness, let $E_2/c\to0$ and let the remaining $n-2$ distinct energies
approach $c$ from below (with $E_n=c$). Then $c\beta_*\to n/(n-2)$.
Any fixed larger value of $\beta_b c$ therefore violates the threshold for
some strictly ordered finite energy vector. The first-contrast construction
in the preceding theorem realizes ordinary cooling with all $G=0$.
\end{proof}

For four states, the spacing-independent exclusion is $\beta_b c\leq2$.
This also explains a limitation of the numerical evidence for the original
conjecture. The March paper's Fig.~2 samples $\beta_bE_i$ between zero and
one, hence $\beta_b c\leq1$. Every triplet-failing sample in that box lies
inside the proved exclusion region, for each of the tested state counts
$n=4,\ldots,7$. Its absence of observed counterexamples in that box is thus
consistent with a stronger restricted theorem, while it cannot justify the
unrestricted inference. The September sampling box with span up to five is
not covered by this corollary.
